\documentclass[10pt,a4paper]{amsart}
\usepackage{amsmath,amssymb,amsthm,graphicx}
% Retained minimal preamble for the Stacks Project chapter more-morphisms.tex (More on Morphisms).
% Verbatim excerpt of the chapter preamble: Xy-pic setup (native preamble lines 14-18) and the native
% theorem declarations with the mathematical macros (native preamble lines 154-217). Copied mechanically;
% nothing added, removed or redefined. The wrapper class and the amsmath/amssymb/amsthm/graphicx base
% packages are supplied by the assembler, so the class/externaldocument part of the native preamble is omitted.
% For commutative diagrams we use Xy-pic
\usepackage[all]{xy}

% We use 2cell for 2-commutative diagrams.
\xyoption{2cell}
% Theorem environments.
%
\theoremstyle{plain}
\newtheorem{theorem}[subsection]{Theorem}
\newtheorem{proposition}[subsection]{Proposition}
\newtheorem{lemma}[subsection]{Lemma}

\theoremstyle{definition}
\newtheorem{definition}[subsection]{Definition}
\newtheorem{example}[subsection]{Example}
\newtheorem{exercise}[subsection]{Exercise}
\newtheorem{situation}[subsection]{Situation}

\theoremstyle{remark}
\newtheorem{remark}[subsection]{Remark}
\newtheorem{remarks}[subsection]{Remarks}

\numberwithin{equation}{subsection}

% Macros
%
\def\lim{\mathop{\mathrm{lim}}\nolimits}
\def\colim{\mathop{\mathrm{colim}}\nolimits}
\def\Spec{\mathop{\mathrm{Spec}}}
\def\Hom{\mathop{\mathrm{Hom}}\nolimits}
\def\Ext{\mathop{\mathrm{Ext}}\nolimits}
\def\SheafHom{\mathop{\mathcal{H}\!\mathit{om}}\nolimits}
\def\SheafExt{\mathop{\mathcal{E}\!\mathit{xt}}\nolimits}
\def\Sch{\mathit{Sch}}
\def\Mor{\mathop{\mathrm{Mor}}\nolimits}
\def\Ob{\mathop{\mathrm{Ob}}\nolimits}
\def\Sh{\mathop{\mathit{Sh}}\nolimits}
\def\NL{\mathop{N\!L}\nolimits}
\def\CH{\mathop{\mathrm{CH}}\nolimits}
\def\proetale{{pro\text{-}\acute{e}tale}}
\def\etale{{\acute{e}tale}}
\def\QCoh{\mathit{QCoh}}
\def\Ker{\mathop{\mathrm{Ker}}}
\def\Im{\mathop{\mathrm{Im}}}
\def\Coker{\mathop{\mathrm{Coker}}}
\def\Coim{\mathop{\mathrm{Coim}}}

% Boxtimes
%
\DeclareMathSymbol{\boxtimes}{\mathbin}{AMSa}{"02}

%
% Macros for moduli stacks/spaces
%
\def\QCohstack{\mathcal{QC}\!\mathit{oh}}
\def\Cohstack{\mathcal{C}\!\mathit{oh}}
\def\Spacesstack{\mathcal{S}\!\mathit{paces}}
\def\Quotfunctor{\mathrm{Quot}}
\def\Hilbfunctor{\mathrm{Hilb}}
\def\Curvesstack{\mathcal{C}\!\mathit{urves}}
\def\Polarizedstack{\mathcal{P}\!\mathit{olarized}}
\def\Complexesstack{\mathcal{C}\!\mathit{omplexes}}
% \Pic is the operator that assigns to X its picard group, usage \Pic(X)
% \Picardstack_{X/B} denotes the Picard stack of X over B
% \Picardfunctor_{X/B} denotes the Picard functor of X over B
\def\Pic{\mathop{\mathrm{Pic}}\nolimits}
\def\Picardstack{\mathcal{P}\!\mathit{ic}}
\def\Picardfunctor{\mathrm{Pic}}
\def\Deformationcategory{\mathcal{D}\!\mathit{ef}}
\newenvironment{reference}{}{}
\newenvironment{slogan}{}{}
\newenvironment{history}{}{}
\usepackage{cleveref}
\crefname{originallemma}{original lemma}{original lemmas}
\Crefname{originallemma}{Original lemma}{Original lemmas}
\makeatletter
\newlabel{algebra-lemma-one-equation}{{algebra-lemma-one-equation}{source}{}{sourcepointer-1}{}}
\newlabel{algebra-lemma-one-equation@cref}{{[originallemma][0][]algebra-lemma-one-equation}{[1][0][]source}{}{}{}}
\makeatother
\crefname{originallemma}{original lemma}{original lemmas}
\Crefname{originallemma}{Original lemma}{Original lemmas}
\makeatletter
\newlabel{divisors-lemma-relative-Cartier}{{divisors-lemma-relative-Cartier}{source}{}{sourcepointer-2}{}}
\newlabel{divisors-lemma-relative-Cartier@cref}{{[originallemma][0][]divisors-lemma-relative-Cartier}{[1][0][]source}{}{}{}}
\makeatother
\crefname{originallemma}{original lemma}{original lemmas}
\Crefname{originallemma}{Original lemma}{Original lemmas}
\makeatletter
\newlabel{lemma-slice-smooth-once}{{lemma-slice-smooth-once}{source}{}{sourcepointer-3}{}}
\newlabel{lemma-slice-smooth-once@cref}{{[originallemma][0][]lemma-slice-smooth-once}{[1][0][]source}{}{}{}}
\makeatother
\crefname{originallemma}{original lemma}{original lemmas}
\Crefname{originallemma}{Original lemma}{Original lemmas}
\makeatletter
\newlabel{morphisms-lemma-closed-point-fibre-locally-finite-type}{{morphisms-lemma-closed-point-fibre-locally-finite-type}{source}{}{sourcepointer-4}{}}
\newlabel{morphisms-lemma-closed-point-fibre-locally-finite-type@cref}{{[originallemma][0][]morphisms-lemma-closed-point-fibre-locally-finite-type}{[1][0][]source}{}{}{}}
\makeatother
\crefname{originallemma}{original lemma}{original lemmas}
\Crefname{originallemma}{Original lemma}{Original lemmas}
\makeatletter
\newlabel{morphisms-lemma-dimension-fibre-at-a-point}{{morphisms-lemma-dimension-fibre-at-a-point}{source}{}{sourcepointer-5}{}}
\newlabel{morphisms-lemma-dimension-fibre-at-a-point@cref}{{[originallemma][0][]morphisms-lemma-dimension-fibre-at-a-point}{[1][0][]source}{}{}{}}
\makeatother
\crefname{originallemma}{original lemma}{original lemmas}
\Crefname{originallemma}{Original lemma}{Original lemmas}
\makeatletter
\newlabel{morphisms-lemma-locally-quasi-finite-rel-dimension-0}{{morphisms-lemma-locally-quasi-finite-rel-dimension-0}{source}{}{sourcepointer-6}{}}
\newlabel{morphisms-lemma-locally-quasi-finite-rel-dimension-0@cref}{{[originallemma][0][]morphisms-lemma-locally-quasi-finite-rel-dimension-0}{[1][0][]source}{}{}{}}
\makeatother
\crefname{originallemma}{original lemma}{original lemmas}
\Crefname{originallemma}{Original lemma}{Original lemmas}
\makeatletter
\newlabel{morphisms-lemma-quasi-finite-points-open}{{morphisms-lemma-quasi-finite-points-open}{source}{}{sourcepointer-7}{}}
\newlabel{morphisms-lemma-quasi-finite-points-open@cref}{{[originallemma][0][]morphisms-lemma-quasi-finite-points-open}{[1][0][]source}{}{}{}}
\makeatother
\crefname{originallemma}{original lemma}{original lemmas}
\Crefname{originallemma}{Original lemma}{Original lemmas}
\makeatletter
\newlabel{schemes-lemma-fibre-topological}{{schemes-lemma-fibre-topological}{source}{}{sourcepointer-8}{}}
\newlabel{schemes-lemma-fibre-topological@cref}{{[originallemma][0][]schemes-lemma-fibre-topological}{[1][0][]source}{}{}{}}
\makeatother
\title{The slice of a smooth morphism by a fibrewise nonzerodivisor is smooth}
\author{The Stacks Project Authors}
\date{stacks-project:more-morphisms.tex:local-snapshot-20261011; GNU FDL 1.2 or later}
\begin{document}
\maketitle
\noindent\textit{Editorial note.} Statement span is lines 10774-10789 of sources/stacks-more-morphisms/pinned-original-source.tex; proof span is lines 10791-10843, the complete author proof (51 body lines).\par\medskip
\noindent\textit{Editorial note.} Navigation stubs: 8 cited-but-not-included labels are supplied as non-load-bearing pointers with their pinned tags where the snapshot tag index has them.\par\medskip
\noindent\textit{Editorial note.} extra\_packages is empty on purpose: the retained preamble already loads xy with options and a second declaration clashes in the assembler wrapper.\par\medskip
\noindent\textit{Editorial note.} Text is retained verbatim, including the author's own typos and placeholders.\par\medskip
\noindent\textit{Allowed standard prerequisites (editorial).}\begin{itemize}
\item Explicit prerequisite boundary. Smooth morphisms and the differential criterion, sheaves of differentials under quotients by a regular element, and nonzerodivisors on fibres; the computation is retained.
\end{itemize}\medskip
\ifdefined\hypertarget\hypertarget{sourcepointer-1}{}\fi
\noindent\textit{Source reference.} Original lemma algebra-lemma-one-equation: sources/stacks-more-morphisms/pinned-original-source.tex, source lines 10843--10843. Cited by the retained statement or proof but not part of the collected spans; supplied only as a navigation pointer. The argument does not depend on its content.\par\medskip
\ifdefined\hypertarget\hypertarget{sourcepointer-2}{}\fi
\noindent\textit{Source reference.} Original lemma divisors-lemma-relative-Cartier: sources/stacks-more-morphisms/pinned-original-source.tex, source lines 10843--10843. Cited by the retained statement or proof but not part of the collected spans; supplied only as a navigation pointer. The argument does not depend on its content.\par\medskip
\ifdefined\hypertarget\hypertarget{sourcepointer-3}{}\fi
\noindent\textit{Source reference.} Original lemma lemma-slice-smooth-once: sources/stacks-more-morphisms/pinned-original-source.tex, source lines 10661--10674. Cited by the retained statement or proof but not part of the collected spans; supplied only as a navigation pointer. The argument does not depend on its content.\par\medskip
\ifdefined\hypertarget\hypertarget{sourcepointer-4}{}\fi
\noindent\textit{Source reference.} Original lemma morphisms-lemma-closed-point-fibre-locally-finite-type: sources/stacks-more-morphisms/pinned-original-source.tex, source lines 10843--10843. Cited by the retained statement or proof but not part of the collected spans; supplied only as a navigation pointer. The argument does not depend on its content.\par\medskip
\ifdefined\hypertarget\hypertarget{sourcepointer-5}{}\fi
\noindent\textit{Source reference.} Original lemma morphisms-lemma-dimension-fibre-at-a-point: sources/stacks-more-morphisms/pinned-original-source.tex, source lines 10843--10843. Cited by the retained statement or proof but not part of the collected spans; supplied only as a navigation pointer. The argument does not depend on its content.\par\medskip
\ifdefined\hypertarget\hypertarget{sourcepointer-6}{}\fi
\noindent\textit{Source reference.} Original lemma morphisms-lemma-locally-quasi-finite-rel-dimension-0: sources/stacks-more-morphisms/pinned-original-source.tex, source lines 10843--10843. Cited by the retained statement or proof but not part of the collected spans; supplied only as a navigation pointer. The argument does not depend on its content.\par\medskip
\ifdefined\hypertarget\hypertarget{sourcepointer-7}{}\fi
\noindent\textit{Source reference.} Original lemma morphisms-lemma-quasi-finite-points-open: sources/stacks-more-morphisms/pinned-original-source.tex, source lines 10843--10843. Cited by the retained statement or proof but not part of the collected spans; supplied only as a navigation pointer. The argument does not depend on its content.\par\medskip
\ifdefined\hypertarget\hypertarget{sourcepointer-8}{}\fi
\noindent\textit{Source reference.} Original lemma schemes-lemma-fibre-topological: sources/stacks-more-morphisms/pinned-original-source.tex, source lines 10843--10843. Cited by the retained statement or proof but not part of the collected spans; supplied only as a navigation pointer. The argument does not depend on its content.\par\medskip
\begin{lemma}
\label{lemma-slice-smooth}
Let $f : X \to S$ be a morphism of schemes.
Let $x \in X$ be a point with image $s \in S$.
Assume
\begin{enumerate}
\item $f$ is smooth at $x$, and
\item $x$ is a closed point of $X_s$ and $\kappa(s) \subset \kappa(x)$
is separable.
\end{enumerate}
Then there exists an immersion $Z \to X$ containing $x$ such that
\begin{enumerate}
\item $Z \to S$ is \'etale, and
\item $Z_s = \{x\}$ set theoretically.
\end{enumerate}
\end{lemma}

\begin{proof}
We may and do replace $S$ by an affine open neighbourhood of $s$.
We may and do replace $X$ by an affine open neighbourhood of $x$
such that $X \to S$ is smooth.
We will prove the lemma for smooth morphisms of affines
by induction on $d = \dim_x(X_s)$.

\medskip\noindent
The case $d = 0$. In this case we show that we may take $Z$ to be
an open neighbourhood of $x$. Namely, if $d = 0$, then $X \to S$
is quasi-finite at $x$, see
Morphisms, Lemma \ref{morphisms-lemma-locally-quasi-finite-rel-dimension-0}.
Hence there exists an affine open neighbourhood $U \subset X$ such
that $U \to S$ is quasi-finite, see
Morphisms, Lemma \ref{morphisms-lemma-quasi-finite-points-open}.
Thus after replacing $X$ by $U$ we see that
$X$ is quasi-finite and smooth over $S$, hence
smooth of relative dimension $0$ over $S$, hence
\'etale over $S$. Moreover, the fibre $X_s$ is a finite
discrete set. Hence after replacing $X$ by a further affine open neighbourhood
of $X$ we see that $f^{-1}(\{s\}) = \{x\}$ (because the topology
on $X_s$ is induced from the topology on $X$, see
Schemes, Lemma \ref{schemes-lemma-fibre-topological}).
This proves the lemma in this case.

\medskip\noindent
Next, assume $d > 0$. Note that because $x$ is a closed point of its
fibre the extension $\kappa(x)/\kappa(s)$ is finite (by the
Hilbert Nullstellensatz, see
Morphisms, Lemma \ref{morphisms-lemma-closed-point-fibre-locally-finite-type}).
Thus we see $\Omega_{\kappa(x)/\kappa(s)} = 0$ as this holds for
algebraic separable field extensions.
Thus we may apply
Lemma \ref{lemma-slice-smooth-once}
to find a diagram
$$
\xymatrix{
D \ar[r] \ar[rrd] & U \ar[r] \ar[rd] & X \ar[d] \\
& & S
}
$$
with $x \in D$. Note that
$\dim_x(D_s) = \dim_x(X_s) - 1$ for example because
$\dim(\mathcal{O}_{D_s, x}) = \dim(\mathcal{O}_{X_s, x}) - 1$ by
Algebra, Lemma \ref{algebra-lemma-one-equation}
(also $D_s \subset X_s$ is effective Cartier, see
Divisors, Lemma \ref{divisors-lemma-relative-Cartier})
and then using
Morphisms, Lemma \ref{morphisms-lemma-dimension-fibre-at-a-point}.
Thus the morphism $D \to S$ is smooth with
$\dim_x(D_s) = \dim_x(X_s) - 1 = d - 1$. By induction hypothesis
we can find an immersion $Z \to D$ as desired, which finishes the proof.
\end{proof}

\end{document}
